\section{引言：长时间运行 Agent 的挑战}

随着 AI Agent 的能力不断增强，开发者开始要求它们承担跨越数小时甚至数天的复杂任务。然而，如何让 Agent 在多个 context window 之间保持一致的进展，仍然是一个开放的工程难题。

\begin{knowledgebox}{什么是 Context Window？}
Context window 是大语言模型在一次推理调用中能够"看到"的全部文本长度。即使像 Claude 这样的前沿模型，其 context window 也是有限的。当一个任务需要多次调用（多个 session）才能完成时，每个新 session 都从零开始——没有上一轮的记忆。这就是长时间运行 Agent 面临的核心问题。
\end{knowledgebox}

Anthropic 工程团队的这篇博客文章提出了一个形象的比喻：想象一个软件项目由多名工程师轮班完成，但每位工程师上班时完全不记得上一班发生了什么。这正是长时间运行 Agent 面临的困境——context window 有限，而大多数复杂项目无法在单个 window 内完成。

为了解决这一问题，Anthropic 开发了一套基于 Claude Agent SDK 的双阶段方案：
\begin{enumerate}
    \item \textbf{Initializer Agent}（初始化 Agent）：在首次运行时搭建环境；
    \item \textbf{Coding Agent}（编码 Agent）：在每个后续 session 中渐进式推进任务，并为下一个 session 留下清晰的工作记录。
\end{enumerate}

\begin{importantbox}{核心思想：像人类工程师一样交接}
这套方案的关键洞察是：让 Agent 像优秀的人类软件工程师那样工作——每次交班前写好交接文档、提交代码、记录进度。通过 \texttt{claude-progress.txt} 文件配合 git history，新的 Agent session 可以快速理解当前项目状态。
\end{importantbox}

\subsection{本章小结}

长时间运行 Agent 的核心挑战在于 context window 的有限性导致跨 session 的状态丢失。Anthropic 的方案采用"初始化 + 渐进式编码"的双阶段架构，通过结构化的环境管理和进度追踪来弥合 session 之间的断裂。


